\documentclass[11pt,xcolor=svgnames]{beamer}

\usepackage[T1]{fontenc}
\usepackage[spanish]{babel}
\usepackage[utf8]{inputenc}

\usepackage{graphicx}
\usepackage{tikz}

\usepackage{charter}
\usepackage{inconsolata}
\usepackage{helvet}

\usetikzlibrary{arrows,shapes,snakes}

\usepackage{colores}

\usetheme{umbc2}
% \usefonttheme{serif}
\usefonttheme{structurebold}
\usecolortheme[named=Gray]{structure}

\title{oFlute: reconocimiento de señales aplicado al aprendizaje de la flauta
  dulce} 

\author{José Tomás Tocino García}

\institute[Universidad de Cádiz]{Universidad de Cádiz}
% \subject{JOSLUCA 4} 

\date[Sept 2011]{Septiembre de 2011}

\defbeamertemplate*{title page}{customized}[1][]
{
  \vbox{}
  \vfill
  \begin{centering}
    \vspace{0.10cm}
    \includegraphics[width=0.215\textwidth]{imagenes/logotipo} 
    \vspace{0.2cm}
    \begin{beamercolorbox}[sep=8pt,center,#1]{title}
      \usebeamerfont{title}\inserttitle\par%
    \end{beamercolorbox}%
    \vskip1em\par
    \begin{beamercolorbox}[sep=0pt,center,#1]{author}
      \usebeamerfont{author}\small\insertauthor\\
      \texttt{<josetomas.tocino@uca.es>}
    \end{beamercolorbox}
    \vspace{0.2cm}
    \begin{beamercolorbox}[sep=0pt,center,#1]{date}
      \usebeamerfont{date}
      {\scriptsize \insertdate }
    \end{beamercolorbox}%\vskip0.5em

    \includegraphics[width=0.12\textwidth]{imagenes/logo_uca} 

    % {\usebeamercolor[fg]{titlegraphic}\inserttitlegraphic\par}
  \end{centering}
  \vfill

}

\setbeamercolor{footline}{fg=white, bg=white}

\begin{document}

% Transparencia de Inicio -> Título
{
  \setbeamertemplate{footline}{}
  \begin{frame}
    \titlepage
  \end{frame}
}
% \logo{\includegraphics[height=2.5cm]{imagenes/logotipo}}
\setbeamertemplate{background}{\includegraphics[width=\paperwidth,height=\paperheight]{imagenes/fondo1.pdf}}
\normalsize

% Transparencia de índice

\begin{frame}{Índice}
  \tableofcontents
\end{frame}


\section{Introducción}

\begin{frame}{Contexto social}
  \begin{center}
    \Large

    Jóvenes en plena simbiosis con las nuevas tecnologías.

    \bigskip
    \pause
    \bigskip

    Las TIC están llegando a los centros educativos.

    \bigskip
    \pause
    \bigskip

    Técnicas docentes basadas en recursos multimedia e informáticos.

  \end{center}
\end{frame}

\begin{frame}{Concepción del proyecto}
  \begin{block}{Idea}
    Hacer un juego educativo.
  \end{block}

  \pause

  \begin{block}{Primera cuestión}
    ¿Sobre qué aspecto educativo? ¿Qué asignatura se beneficia?
  \end{block}

  \vspace{0.5cm}
  \pause

  \begin{center}
    \LARGE Música \\[0.5cm]

    \pause

    \LARGE Aprendizaje de la flauta dulce
  \end{center}
\end{frame}

% \begin{frame}{Contexto}
%   \begin{block}{Wikipedia}
%     \noindent \LaTeX{} es un \textbf{lenguaje de marcado} para
%     documentos, y un sistema de preparación de documentos, formado
%     por un gran conjunto de macros de \TeX{}
%   \end{block}

%   \pause

%   \begin{block}{Más normalito...}
%     \noindent \LaTeX{} es un ``lenguaje de programación'' que nos
%     permite editar documentos con muchas posibilidades, como fórmulas
%     matemáticas, tablas y muchos más.
%   \end{block}

%   \pause

%   \begin{block}{\textbf{¡IMPORTANTE!}}
%     \noindent Los documentos \LaTeX{} usan extensión \texttt{.tex}, y
%     son ficheros de texto plano, que luego se compilan, del tipo
%     \textbf{WYSINWYG}.
%   \end{block}

% \end{frame}

\end{document}

